\section{Proposed Method}

MemeSonic is organized into two phases (Figure~\ref{fig:pipeline}): (1) Incongruity-Aware
Multimodal Understanding, which produces a structured affective representation of the meme;
and (2) Meme-to-Speech Generation, which translates this representation into expressively appropriate audio.

\subsection{Phase 1: Incongruity-Aware Multimodal Understanding}

\subsubsection{Explicit Affect Decomposition}
We encode the meme image $v$ with the CLIP image encoder $\phi_v$ and text $t$ with the CLIP text
encoder $\phi_t$, both frozen. Each encoder output is projected onto a set of $C$ emotion prototype
vectors $\{\mathbf{e}_c\}_{c=1}^C$ derived from CLIP text encodings of emotion label descriptions
(e.g., ``a happy and joyful image''):
\begin{equation}
  p^{(\cdot)}_c = \frac{\exp\!\bigl(\phi_{(\cdot)}(\cdot)^\top \mathbf{e}_c / \tau\bigr)}
                       {\sum_{j=1}^C \exp\!\bigl(\phi_{(\cdot)}(\cdot)^\top \mathbf{e}_j / \tau\bigr)}
\end{equation}
yielding emotion distributions $p_v, p_t \in \Delta^C$ for image and text respectively, where
$\tau$ is a temperature parameter. The incongruity score is:
\begin{equation}
  \delta(v,t) = \tfrac{1}{2}\,\|p_v - p_t\|_1 \;\in [0,1]
  \label{eq:incong}
\end{equation}
A value near 0 indicates emotional alignment (redundancy); a value near 1 indicates strong conflict
(synergy/incongruity). This scalar is differentiable with respect to the projection parameters and
is used throughout both phases.

\subsubsection{MMOE Fusion with Incongruity Routing}
\label{sec:mmoe}
Rather than applying a fixed fusion strategy, we route each sample to one of three expert fusion
modules based on its incongruity score $\delta$:

\textbf{Expert R (Redundancy)} handles samples where image and text are emotionally aligned
($\delta$ low). It concatenates the projected image and text embeddings:
$z_R = [\phi_v(v);\, \phi_t(t)]$, leveraging shared information.

\textbf{Expert U (Uniqueness)} handles samples where one modality dominates. It applies late fusion
via a learned weighted combination:
$z_U = \alpha \phi_v(v) + (1-\alpha) \phi_t(t)$, preserving modality-specific contributions.

\textbf{Expert S (Synergy)} handles high-incongruity samples. It applies Incongruity Fusion:
\begin{equation}
  z_S = \text{MLP}\bigl([\phi_v(v);\, \phi_t(t);\, p_v - p_t;\, \delta]\bigr)
\end{equation}
directly encoding the conflict vector $p_v - p_t$ and scalar $\delta$ into the representation.

The router is a single linear layer conditioned on $[\phi_v(v);\, \phi_t(t);\, \delta]$:
\begin{equation}
  \mathbf{r} = \text{softmax}\bigl(W_r [\phi_v(v);\, \phi_t(t);\, \delta] + b_r\bigr)
\end{equation}
The final fused representation is a soft mixture:
\begin{equation}
  z = r_R \cdot z_R + r_U \cdot z_U + r_S \cdot z_S
\end{equation}
followed by a task-specific linear classifier. During training, we use a two-objective loss:
$\mathcal{L} = \mathcal{L}_\text{cls}(z, y) + \lambda \mathcal{L}_\text{gate}(\mathbf{r}, \mathbf{r}^*)$,
where $\mathbf{r}^*$ are soft gate labels derived from the per-sample performance of the trained Concat
and Late Fusion experts.

\subsection{Phase 2: Meme-to-Speech Generation}

\subsubsection{Script Generation}
Given the meme image and text, we use GPT-4.1-mini with a multimodal prompt that provides the
image, the detected text, and the predicted sentiment/intention labels to generate a voiceover script
$s$. The prompt instructs the model to write a short narration that reflects the \emph{pragmatic intent}
of the meme—including irony and sarcasm—rather than a literal description.

\subsubsection{Incongruity-to-Prosody Adapter (Closed-Source Path)}
For the closed-source TTS path (ElevenLabs), the fused representation $z$ and incongruity score $\delta$
are passed through a fine-tuned LLM adapter that generates a natural-language prosody prompt:
$\text{prompt}_\text{prosody} = \text{LLM\_adapter}(z, \delta, s)$, e.g., ``Read with deadpan irony,
slightly elevated pitch, and a deliberate pause before the punchline.''
The prosody prompt, script, and a fixed voice identity are then sent to ElevenLabs to synthesize audio.
The voice settings (stability, style) are additionally modulated as a function of $\delta$: high
incongruity leads to lower stability (more expressive variance) and higher style (more dramatic delivery).

\subsubsection{Embedding-Guided Open-Source Path}
For the open-source path (StyleTTS2), the fused latent $z \in \mathbb{R}^d$ is passed through a
lightweight adapter (2-layer MLP) that predicts StyleTTS2 style vector $\hat{s} \in \mathbb{R}^{d_s}$:
$\hat{s} = \text{Adapter}(z)$.
The adapter is trained on pseudo-GT audio generated by the closed-source path, using a cosine
embedding loss: $\mathcal{L}_\text{adapter} = 1 - \cos(\text{StyleEmb}(\Psi(\hat{s}, s)),\, \text{StyleEmb}(a^*))$,
where $a^*$ is the pseudo-GT audio. At inference, StyleTTS2 synthesizes audio conditioned on $\hat{s}$.

\subsection{Tri-Modal Alignment: Linear Projector for Audio Integration}
\label{sec:projector}

To diagnose why naive tri-modal late fusion fails and whether audio carries independent affective
information, we train a linear projector $W \in \mathbb{R}^{D_a \times D_{ib}}$ that maps
Emotion2Vec audio embeddings~\citep{ma2023emotion2vec} ($D_a=1024$) into ImageBind's shared
embedding space ($D_{ib}=1024$):
\begin{equation}
  \hat{a} = \text{normalize}\bigl(W\, \phi_a(a) + b\bigr)
\end{equation}
The projector is supervised with MSE loss against the ImageBind text prototype of the sample's
ground-truth sentiment class, using all 9 voice-setting variants as data augmentation:
\begin{equation}
  \mathcal{L}_\text{proj} = \frac{1}{N} \sum_{i=1}^N
    \bigl\|\hat{a}_i - \text{normalize}(\mathbf{e}_{y^{(s)}_i})\bigr\|_2^2
\end{equation}
After projection, tri-modal fusion is performed by averaging $\phi_v(v)$, $\phi_t(t)$, and $\hat{a}$
in the shared ImageBind space, followed by cosine-similarity classification against text prototypes.
This design decouples the diagnosis of (1) modality-space incompatibility from
(2) the intrinsic informativeness of audio embeddings.
